\section{The Evidence for a Life}

\subsection{What Tertullian says, and does not say, about himself}

The manuscripts of his works transmit the full name Quintus Septimius Florens
Tertullianus. Only once does the author name himself: the treatise \work{On the
Veiling of Virgins} closes with a third-person subscription wishing peace to
willing readers ``as likewise to Septimius Tertullianus, whose this tractate
is'' (\work{On the Veiling of Virgins} 17). Everything else must be inferred
from some thirty-one surviving Latin treatises, none of them autobiographical.

The inferences are nonetheless substantial. He wrote a supple, aggressive,
learned Latin saturated with the techniques of the rhetorical schools, and he
handled philosophical, medical, legal, and antiquarian material with the ease
of an expensively educated man. He knew Greek well enough to compose in it: ``this point has already received
a fuller discussion from us in Greek,'' he remarks of a baptismal question
(\work{On Baptism} 15)---one of several notices of Greek works, all now lost. He was married; the two books \work{To
His Wife} address her directly and ask that ``after our departure,'' she
``renounce nuptials'' (\work{To His Wife} 1.1). He was a convert who remembered
being otherwise: he exhorts ``that repentance, O sinner, like myself,'' be
embraced ``as a shipwrecked man the protection of some plank'' (\work{On
Repentance} 4), and his works assume intimate former acquaintance with the
games, the theatre, and pagan civic religion. Most strikingly, he twice speaks
of himself as one of the laity. Arguing that in extremity every Christian may
baptize and offer, he asks, ``Are not even we laics priests? \ldots\ where
three are, a church is, albeit they be laics'' (\work{On Exhortation to
Chastity} 7); and in \work{On Monogamy} he mocks how ``we,'' the laity, exalt
ourselves against the clergy when it suits us (\work{On Monogamy} 12).

\subsection{The received portrait: Jerome and Eusebius}

The earliest connected biography is a paragraph written about 392--393. Jerome
of Stridon, in his catalogue \work{On Illustrious Men}, makes Tertullian ``the
presbyter, now regarded as chief of the Latin writers,'' a Carthaginian, ``the
son of a proconsul or centurion, a man of keen and vigorous character''; he
flourished under Severus and Caracalla; he ``was presbyter of the church until
middle life, afterwards driven by the envy and abuse of the clergy of the Roman
church, he lapsed to the doctrine of Montanus''; in old age he wrote many works
now lost, and ``he is said to have lived to a decrepit old age'' (\work{On
Illustrious Men} 53). Jerome also preserves the most famous item of ancient
reception: an old man of Concordia who had known Cyprian's secretary told him
that Cyprian read Tertullian daily and would ask for the book with the words
``Give me the master'' (\textit{Da magistrum}) (ibid.).

Almost a century before Jerome, Eusebius of Caesarea had cited the
\work{Apology}---which he knew in a Greek translation---and introduced its
author as ``a man well versed in the laws of the Romans, and in other respects
of high repute, and one of those especially distinguished in Rome''
(\work{Church History} 2.2.4; cf.\ 5.5, where Tertullian is ``a trustworthy
witness''). Eusebius thus places Tertullian, at least by reputation, in a Roman
and legal milieu.

\subsection{The modern audit of the portrait}

Twentieth-century scholarship subjected this portrait to a severe audit, most
influentially Timothy Barnes's \work{Tertullian: A Historical and Literary
Study} (1971). Barnes argued that nearly every biographical detail not found in
Tertullian's own works dissolves on inspection: ``Tertullian's father was not a
centurion; there never was a \textit{centurio proconsularis}''---the office
rests on a corrupted text of Jerome; the presbyterate contradicts Tertullian's
own lay self-description; the identification of the Christian writer with the
jurist Tertullianus excerpted in the \work{Digest} is unfounded; Eusebius's
Roman legal expert is an inference from the \work{Apology}'s forensic manner;
and the tale that envy of the Roman clergy drove him to Montanism explains
nothing in the African evidence. On Barnes's minimal reconstruction, Tertullian
was a lay, married, classically educated Carthaginian who wrote intensively
between about 196 and 212 and died not long afterwards. Reviewers accepted much
of the demolition while doubting the compression: W.\,H.\,C.\ Frend, among
others, held that Jerome ``is unlikely to have been completely mistaken'' on
every point and that the traditional longer life---birth around 155--160,
death, on Jerome's report of great old age, perhaps as late as the 220s---could
not be excluded (Frend 1974). No consensus has replaced the ranges: birth
somewhere in c.~155--170, death after 212 at an unknown date.

\witnesslimit{Every element of the customary summary ``born c.~155, son of a
centurion, lawyer in Rome, priest of Carthage, died c.~220'' is either a late
report, a disputed reading, or a modern inference. This study therefore treats
only the self-witness and the dated works as controlling, and reports the rest
as tradition under review.}
